\documentclass[10pt]{article}

\usepackage[utf8]{inputenc}

\usepackage[T1]{fontenc}

\usepackage{amsmath}

\usepackage{amsfonts}

\usepackage{amssymb}

\usepackage[version=4]{mhchem}

\usepackage{stmaryrd}

\usepackage{graphicx}

\usepackage[export]{adjustbox}

\graphicspath{ {./images/} }

\usepackage{bbold}



\begin{document}

При этом $x=x_{+}-x_{-}$, и эта формула дает «минимальное" представление $x$ в виде разности положительных элементов, т. е. если $x=y-z$, где $y, z \geqslant 0$, то $x_{+} \leqslant y$, $x_{-} \leqslant z$. Миниэдральный конус является воспроизводящим. Элемент $|x|=x_{+}+x_{-}$наз. модулем элемента $x$. В пространстве $C[a, b]$ с естественным упорядочением положительный конус миниэдрален, положительная часть любой функции $x(t)$ из $C$ получается из $x(t)$ заменой ее отрицательных значений нулем, а модуль есть функция $|x(t)|$. В в. р. всякое конечное множество элементов имеет обе грани. Модуль элемента в. р. обладает многими свойствами абсолютной величины действительного числа.



В. р. наз. д и с т р и б у т и в н о й, если для произвольного множества ее элементов $\left\{x_{\alpha}\right\}$, у к-рого существует $\sup x_{\alpha}$, при любом $y \in X$ справедлива формула. $y \wedge$ $\wedge \sup x_{\alpha}=\sup \left\{y \wedge x_{\alpha}\right\}$. Тогда верна и двойственная формула. $y \vee \inf x_{\alpha}=$ inf $\left\{y \vee x_{\alpha}\right\}$.



\includegraphics[max width=\textwidth, center]{2023_07_08_b89736afc0f7353b3e37g-1}

жительных элементов: если $x=y+z$, где $y, z \geqslant 0$, и одновременно $x=x_{1}+\ldots+x_{n}$, где все $x_{i} \geqslant 0$, то каждый $x_{i}$ можно представить в виде $x_{i}=y_{i}+z_{i}$ так, что все $y_{i}, z_{i} \geqslant 0$ и что



\[

y=y_{1}+\ldots+y_{n}, z=z_{1}+\ldots+z_{n} .

\]



Два элемента $x, y$ в. р. наз. д и 3 ъ ю н к т ны м и $(x d y)$, если $|x| \wedge|y|=0$. Два множества $A, B$ наз. д и з тю н к тными, если $a d b$ для любых $a \in A, b \in B$. В пространстве $C[a, b]$ дизъюнктность $x d y$ означает, что $x(t) y(t) \equiv 0$. Положительный элемент $e$ наз. с л а б о й е дини пей (е динйей в смысле Ф р ейд е н т а л я), если 0 - единственный элемент, дизъюнктный с $e$. В $C[a, b]$ слабой единицей является любая функция, к-рая больше 0 на всюду плотном множестве. Если же әлемент $e$ таков, что для любого $x$ существует $\lambda$, при к-ром $|x| \leqslant \lambda e$, то $e$ наз. с и ль ь о й е д и н иц е й, а $X$ с сильной единицей наз. в. р. о г р а в ич е н ны х э л е м е н т о в. В $C[a, b]$ сильная единица - любая функция, для к-рой $\min x(t)>0$. Если в архимедовой в. p. $X$ с сильной единицей $e$ положить $\|x\|=\min \{\lambda:|x| \leqslant \lambda e\}$, то $X$ становится н о р м и р ов а нн о й $\mathrm{p}$ е пе е т ко й.



На плоскости любой конус, кроме одномерного (т. е. луча), миниэдрален. Но в пространствах с большим числом измерений среди замкнутых конусов много не миниэдральных, напр. таковы все «круглые» конусы в $\mathbb{R}^{3}$. Для того чтобы конус (с вершиной в нуле) в $n$-мерном архимедовом у. в. п был миниәдральным, необходимо и достаточно, чтобы он был натянут на $(n-1)$-мерный симплекс с линейно независимыми вершинами. Всякая архимедова $n$-мерная в. р. изоморфна пространству $\mathbb{R} n$ с покоординатным упорядочением.



К-пространства (К.-П.), п р о с т р а н с т в а К а нт о р о в и ч а, суть (o)-полные в. р. Это- основной класс П. п., они всегда архимедовы. (o)-сходимость в К.-П. описывается с помощью верхнего и нижнего пределов, именно, для огранич. последовательности $\left\{x_{n}\right\}$



\[

\overline{\lim } x_{n}=\inf _{n} \sup _{m \geqslant n} x_{m}, \underline{\lim } x_{n}=\sup _{n} \inf _{m \geqslant n} x_{m}

\]



и тогда $x_{n} \stackrel{(0)}{\longrightarrow} x$ означает, что $x=\varlimsup \lim x_{n}=\lim x_{n}$. Пусть $X$ есть К.-П. Для любого множества $E \subset X$ его дизъюнктным дополнением наз. множество $E^{d}=\{x \in X: x d E\}$. Всякое множество, к-рое является дизъюнктным дополнением (к какому-либо множеству), наз. п о л о с о й. Для любого множества $E$ существует наименьшая полоса, содержащая $E$, именно $E^{d d}$; она наз. п о л о с о й, п о ро ж д е н о й мно же с т в о м $E$. Если само $E$ есть полоса, то $E d d=E$. Полоса, порожденная одноэлементным множеством, наз. г л а в н о й. Понятие полосы вводится и в любой в. р., однако в К.-п оно играет особую роль, поскольку справедлива т е о р е м а о п р о е к т и р о в а и и на полосу: если $E-$ полоса в $X$, то для любого $x \in X$ существует единственное разложение $x=y+z$, где $y \in E, z \in E^{d}$. Определенный при этом линейный оператор $y=\operatorname{Pr}_{E} x \quad$ наз. п р о е к т о р о м н а п о л о с у $E$. Если задана произвольная совокупность попарно дизъюнктных полос $E_{\alpha}$, полная в том смысле, что 0 - единственный элемент из $X$, дизъюнктный всем $E_{\alpha}$, то любой $x \in X_{+}$представим в виде $x=$ $=\sup x_{\alpha}$, где $x_{\alpha}=\operatorname{Pr}_{E \alpha} x$. Всякий $l$-идеал $Y \subset X$ также является К.-П. Однако если $x_{n} \in Y$ и $x_{n} \stackrel{(0)}{\longrightarrow} x$ в $X$, то это соотношение верно и в $Y$ только в том случае, когда последовательность $\left\{x_{n}\right\}$ ограничена в $\boldsymbol{Y}$.



Примером К.-п. служит пространство $S$ всех действительных почти всюду конечных измеримых функций на $[0,1]$, в к-ром эквивалентные функции отождествляются. Функция $x \in S$ считается пөложительной, если $x(t) \geqslant$ $\geqslant 0$ почти всюду. Если $A=\left\{x_{n}\right\}$ - счетное ограниченное сверху подмножество из $S$ (ограниченность сверху означает, что существует такая $y \in S$, что $x_{n}(t) \leqslant y(t)$ почти всюду для любого $n)$, то функция $x(t)=\sup x_{n}(t)$ и будет верхней гранью множества $A$, т. е. sup $A$ вычисляется поточечно. Однако для несчетных множеств вычисление граней таким же способом уже невозможно, и существование у несчетных множеств ограниченных множеств в $S$ устанавливается сложнее. (o)-сходимость в $S$ означает сходимость почти всюду. Все пространства $L_{p}=[0,1], p>0$, являются $l$-идеалами в $S$, и потому они тоже являются К.-П.



Важную роль играет теорема Рисса - Канторовича о том, что множество всех по р я д к о в о о г р а н ич е н ны х п е р а т о р в (т. е. линейных операторов, переводящих ограниченные по упорядочению множества в такие же) из в. р. в К.-П. при естественном порядке $(A \leqslant B$ означает, что $A x \leqslant B x$ при всех $x \supseteqq 0)$ само является К.-П. Теория К.-п. нашла приложения в выпуклом анализе и теории экстремальных задач. Многие результаты здесь основываются на теореме Хана - Банаха - Канторовича о продолжении линейных операторов со значениями в К.-П.



К.-П. наз. р а с ш и р е н н ы м, если всякое множество его попарно дизъюнктных элементов ограничено. В расширенном К.-п. всегда существует слабая единица. Для любого К.-П. $X$ существует единственное (с точностью до изоморфизма) расширенное К.-П. $Y$, в к-рое $X$ погружается как $l$-идеал, а полоса в $Y$, порожденная $X$, совпадает с $Y$. Такое $Y$ наз. м а к с и м а л ь н ы м р а с ши р е в и е м К. - П. Пространство $S[0,1]-$ максимальное расширение для всех пространств $L_{p}[0$, 1]. Понятие расширенного К.-П. играет существенную роль во всей теории П. п., в частности при представлении К. п. с помощью функций.



С в. р. и К.-П. тесно связано понятие решеточнонормированного простравства - векторного пространства, каждому элементу к-рого сопоставлена его обобщенная норма, являющаяся элементом фикспрованной в. р. и удовлетворяющая обычным для нормы аксиомам, в к-рых знак неравенства понимается в смысле порядка в указанной в. р. Такие пространства используются в теории функциональных уравнений (теоремы существования; методы приближенного решения; метод Ньютона - Канторовича; монотовные процессы последовательных приближений и т. д.).



Топологические полуупорядоченные пространства. В функциональном анализе используются также у . в. п., в к-рых одновременно определена еще нек-рая топология, согласованная с порядком. Простейший и важнейший пример такого пространства - банахова peшетка. Обобщением понятия банаховой решетки служит ло к а ль н о в ы п у клая р е ше т к а. Важный класс банаховых К.-П. составляют КВп рост ра н с тва (КВ-П.), Канторовита-





\end{document}